\documentclass[10pt,a4paper]{amsart}
\usepackage[margin=19mm]{geometry}
\usepackage{amsmath,amssymb,mathtools}
\usepackage[scr=rsfs,cal=euler]{mathalfa}
\usepackage{xfrac}
\usepackage[unicode,hidelinks,bookmarks=false]{hyperref}
\newcommand{\ie}{i.e.\ }
\newcommand{\cf}{cf.\ }
\newcommand{\resp}{respectively\ }
\newcommand{\Subsection}{\subsection}
\providecommand{\nobreakdash}{\nobreak\hbox{-}}
\setlength{\emergencystretch}{2em}
\multlinegap0pt
\usepackage{amssymb}
\usepackage{mathtools}
\usepackage[shortlabels]{enumitem}
\usepackage{xcolor}
\newtheorem{theorem}{Theorem}
\newcommand{\Ree}{\operatorname{Re}}
\newcommand{\ZF}{Z_F}
\title{\bfseries The Right Edge of the Zero Set of the Fibonacci Zeta Function}
\author{Marco Mantovanelli}
\date{Source: arXiv:2609.04993v1 (CC BY 4.0)}
\begin{document}
\maketitle

\noindent\emph{Source and licence.} \bfseries The Right Edge of the Zero Set of the Fibonacci Zeta Function. Version arXiv:2609.04993v1 (CC BY 4.0), distributed by arXiv under the Creative Commons Attribution 4.0 International licence, http://creativecommons.org/licenses/by/4.0/, as recorded in the arXiv licence metadata for this exact version and shown on the abstract page https://arxiv.org/abs/2609.04993v1. Full licence notice: \texttt{licenses/arxiv-2609.04993-CC-BY-4.0.txt}.\par\smallskip

\noindent\textit{Editorial scope.} Counted unit: the article's theorem thm:completion (source0.tex lines 1125--1150). The article's complete original proof of it is delivered with it (source0.tex lines 1152--1235, one \texttt{proof} environment of 84 lines). No mathematics is written, repaired or re-derived: the statement and the proof are the article's own lines, in the article's own notation, and no retained line is substituted or re-flowed.  Retained uncounted as the source-local context the proof needs: lines 1093--1097; lines 1099--1103; lines 1119--1123.  Nothing was omitted from any retained span, so the package carries no omission record.  No retained line contains a citation, so the wrapper carries no bibliography and no bibliography entry.  No retained line contains a figure environment, an image-inclusion command or drawing code, so no image is delivered and the package declares no asset dependency.  Editorial formatting only, in addition: an AMS \texttt{amsart} preamble that restates the article's own declarations and macros used by the retained lines, namely the environments \texttt{theorem} on separate counters, together with the standard packages those lines need.  The retained environments print in this document as Theorem 1 in order of appearance, whereas the article prints the same items with the numbering of its own full document.  The complete native source of the article as distributed in the arXiv e-print for this exact version is bundled unchanged as \texttt{sources/209404/source0.tex}.
\begin{equation}\label{eq:global-continuation}
\ZF(s)
=5^{s/2}\sum_{k\ge0}\binom{-s}{k}
\frac{1}{\varphi^{s+2k}+(-1)^{k+1}},
\end{equation}

\begin{equation}\label{eq:pole-lattice}
p_{k,m}
=-2k+\frac{(2m+k)\pi i}{\log\varphi},
\qquad k\ge0,\quad m\in\mathbb Z.
\end{equation}

\begin{equation}\label{eq:completion-def}
P_F(s):=(\varphi^{-s};q)_\infty,
\qquad
\Xi_F(s):=P_F(s)\ZF(s).
\end{equation}

\begin{theorem}[Global completion]\label{thm:completion}
The function $\Xi_F$ extends to an entire function.  Its zeros, counted with multiplicity, are exactly the zeros of the meromorphic function $\ZF$: none of the poles in \eqref{eq:pole-lattice} becomes a zero after cancellation.  More precisely,
\begin{equation}\label{eq:completion-pole-value}
\Xi_F(p_{k,m})
=q^{-k(k+1)/2}(q;q)_k(q;q)_\infty
\,5^{p_{k,m}/2}\binom{-p_{k,m}}{k}\neq0.
\end{equation}
If
\[
M_\Xi(r)=\max_{|s|\le r}|\Xi_F(s)|,
\]
then
\begin{equation}\label{eq:completion-type}
\limsup_{r\to\infty}\frac{\log M_\Xi(r)}{r^2}
=\frac{\log\varphi}{4}.
\end{equation}
Thus $\Xi_F$ has exact order $2$ and exact type $\log\varphi/4$.

Let $n_\infty(R)$ be the number of poles of $\ZF$ in $|s|\le R$, and let $n_0(R)$ be the number of zeros there, both counted with multiplicity.  Then
\begin{equation}\label{eq:global-counting}
n_\infty(R)
=\frac{\log\varphi}{8}R^2+O(R),
\qquad
n_0(R)=O(R^2).
\end{equation}
\end{theorem}

\begin{proof}
The $k$th denominator in \eqref{eq:global-continuation} factors as
\[
\varphi^{s+2k}+(-1)^{k+1}
=\varphi^{s+2k}\bigl(1-\varphi^{-s}q^k\bigr).
\]
Consequently the zeros of $P_F$ are precisely the points \eqref{eq:pole-lattice}, and they are simple.  At $p=p_{k,m}$ one has $\varphi^{-p}q^k=1$, and therefore
\begin{align*}
P_F'(p)
&=\Lambda\prod_{j\ne k}(1-q^{j-k})\\
&=\Lambda(-1)^kq^{-k(k+1)/2}(q;q)_k(q;q)_\infty.
\end{align*}
On the other hand, the residue of \eqref{eq:global-continuation} at the same point is
\[
\operatorname*{Res}_{s=p}\ZF(s)
=\frac{5^{p/2}\binom{-p}{k}}
{\Lambda(-1)^k}.
\]
Multiplying these two expressions proves \eqref{eq:completion-pole-value}.  The right-hand side is nonzero: $(q;q)_\infty\ne0$ since $|q|<1$, and $\Ree(-p_{k,m})=2k$ prevents $\binom{-p_{k,m}}{k}$ from vanishing.  Hence every singularity is removable in \eqref{eq:completion-def}, and the cancellation introduces no new zero.

For the growth estimate, write
\[
P_{F,k}(s)=\prod_{\substack{j\ge0\\j\ne k}}
(1-\varphi^{-s}q^j).
\]
For every fixed $k$, the product $P_{F,k}$ converges locally uniformly and defines an entire function.  Moreover, for $|s|\le r$, $r\ge1$,
\[
|P_{F,k}(s)|
\le Q(r):=\prod_{j\ge0}(1+\varphi^{r-2j}).
\]
With $x=\varphi^{-2}$ one also has
\[
\sum_{k\ge0}\left|\binom{-s}{k}\right|x^k
\le (1-x)^{-r}.
\]
Consequently
\[
\sum_{k\ge0}\left|\binom{-s}{k}\right|\varphi^{-2k}|P_{F,k}(s)|
\le Q(r)(1-\varphi^{-2})^{-r},
\]
so the series below converges absolutely and locally uniformly on $\mathbb C$ by the Weierstrass test.  Off the pole lattice direct cancellation in \eqref{eq:global-continuation} gives
\begin{equation}\label{eq:Xi-series}
\Xi_F(s)
=5^{s/2}\varphi^{-s}
\sum_{k\ge0}\binom{-s}{k}\varphi^{-2k}P_{F,k}(s),
\end{equation}
and since the right-hand side is entire, the identity holds everywhere by continuation.

Splitting the logarithm of $Q(r)$ at $j=r/2$ gives
\begin{equation}\label{eq:qproduct-growth}
\sum_{j\ge0}\log(1+\varphi^{r-2j})
\le \frac{\Lambda}{4}r^2+O(r).
\end{equation}
The remaining prefactor in \eqref{eq:Xi-series}, as well as the factor $(1-\varphi^{-2})^{-r}$ already displayed above, is only exponential in $r$.  Thus
\[
\log M_\Xi(r)\le \frac{\Lambda}{4}r^2+O(r).
\]

For the reverse bound, use the real pole subsequence $p_{2N,-N}=-4N$.  Formula \eqref{eq:completion-pole-value} yields
\[
|\Xi_F(-4N)|
=|q|^{-N(2N+1)}
|(q;q)_{2N}(q;q)_\infty|
\,5^{-2N}\binom{4N}{2N}.
\]
Since $(q;q)_{2N}\to(q;q)_\infty\ne0$, Stirling's formula gives
\[
\log|\Xi_F(-4N)|
=4\Lambda N^2+O(N)
=\frac{\Lambda}{4}(4N)^2+O(N),
\]
which proves \eqref{eq:completion-type}.

Finally, for fixed $k$ the poles in \eqref{eq:pole-lattice} have vertical spacing $2\pi/\Lambda$.  Hence
\begin{align*}
n_\infty(R)
&=\sum_{0\le k\le R/2}
\left(\frac{\Lambda}{\pi}\sqrt{R^2-4k^2}+O(1)\right)\\
&=\frac{\Lambda}{\pi}
\int_0^{R/2}\sqrt{R^2-4x^2}\,dx+O(R)\\
&=\frac{\Lambda}{8}R^2+O(R).
\end{align*}
The bound $n_0(R)=O(R^2)$ follows from Jensen's formula applied to the entire function $\Xi_F$ and the order-two bound above.
\end{proof}

\end{document}
